\documentclass[12pt,a4paper]{article}

% ─── Fonts & Vietnamese ───────────────────────────────────────
\usepackage[utf8]{inputenc}
\usepackage[T5]{fontenc}
\usepackage[vietnamese,english]{babel}
\usepackage{times}

% ─── Packages ──────────────────────────────────────────────────
\usepackage{amsmath,amssymb}
\usepackage{booktabs}
\usepackage{graphicx}
\usepackage{hyperref}
\usepackage[margin=2.5cm]{geometry}
\usepackage{enumitem}
\usepackage{fancyhdr}
\usepackage{setspace}
\usepackage{xcolor}
\usepackage{tcolorbox}
\onehalfspacing

% ─── Header/Footer ────────────────────────────────────────────
\pagestyle{fancy}
\fancyhf{}
\lhead{Báo cáo Tuần 10}
\rhead{[TODO: Tên SV] -- [TODO: MSSV]}
\cfoot{\thepage}

\begin{document}

% ─── Title Page ───────────────────────────────────────────────
\begin{center}
  {\LARGE\bfseries Báo cáo Tuần 10}\\[4pt]
  {\large Báo cáo Thực tập Cơ sở}\\[12pt]
  [TODO: Tên sinh viên]\\
  Lớp: [TODO] \quad MSSV: [TODO]\\
  Giảng viên hướng dẫn: [TODO: Học hàm/học vị. Tên GVHD]\\[4pt]
  Ngày [TODO] tháng 4 năm 2026
\end{center}

\vspace{1em}

% ══════════════════════════════════════════════════════════════
\section{Tổng quan Tuần 10 và Chuyển hướng Nghiên cứu}

\subsection{Kết quả đạt được Tuần 9 (Tóm tắt)}

Tuần 9 đã hoàn thành toàn bộ pipeline debug và cải tiến kiến trúc cho bài toán \emph{reconstruction} (khôi phục) trường mực nước Vịnh Bắc Bộ:

\begin{itemize}[nosep]
  \item Fix 4 lỗi hệ thống: C-property solver, tidal forcing, ocean sensors, ocean-only labels.
  \item Fourier DeepONet: Rel L2 giảm từ 93,6\% $\to$ \textbf{14,8\%} trên test set stratified.
  \item Inference nhanh gấp $\sim$2.300$\times$ so với Godunov HLL solver.
\end{itemize}

\subsection{Chuyển hướng chiến lược: Từ Reconstruction sang Causal Forecasting}

Sau khi phân tích lại toàn bộ câu chuyện nghiên cứu, chúng tôi nhận ra rằng kết quả \emph{reconstruction} (non-causal) tuy ấn tượng nhưng \textbf{chưa phản ánh đúng bài toán thực tế}. Trong thực tế, ta không bao giờ có sẵn chuỗi thời gian tương lai tại các trạm---ta chỉ có dữ liệu quá khứ.

\begin{tcolorbox}[colback=blue!5, colframe=blue!50!black, title={Tái định vị bài toán trung tâm}]
\textbf{Bài toán cũ (reconstruction):} Biết toàn bộ 168h sensor → khôi phục field tại mọi $t$.\\[4pt]
\textbf{Bài toán mới (causal forecasting):} Chỉ biết sensor từ $t=0$ đến $T_\mathrm{obs}$ → \textbf{dự báo} field tại $t > T_\mathrm{obs}$.
\end{tcolorbox}

Đây là bước chuyển đổi từ ``chứng minh khái niệm'' sang ``giải bài toán có ý nghĩa ứng dụng thực sự''.

% ══════════════════════════════════════════════════════════════
\section{Vì sao Sparse Boundary Sensing vẫn là bài toán quan trọng?}

Một câu hỏi tự nhiên: \emph{với vệ tinh altimetry hiện đại (Jason-3, Sentinel-6, SWOT), liệu bài toán ``ít sensor ven bờ'' còn có ý nghĩa?} Câu trả lời dứt khoát là \textbf{có}, với 4 lý do khoa học vững chắc:

\subsection{Vùng ven bờ vẫn là thách thức với altimetry}

Radar altimetry đạt độ chính xác cao ở đại dương mở, nhưng vùng nearshore (cách bờ $<$20\,km) vẫn bị ảnh hưởng nghiêm trọng bởi:
\begin{itemize}[nosep]
  \item \textbf{Land contamination:} Tín hiệu radar bị nhiễu bởi phản xạ từ đất liền nằm trong footprint~[1].
  \item \textbf{Sóng ngắn và động lực ven bờ phức tạp:} Waveform retracking khó hơn nhiều so với open ocean.
  \item \textbf{Correction errors:} Các mô hình hiệu chỉnh khí quyển (wet troposphere, ionosphere) kém chính xác hơn gần bờ.
\end{itemize}

Các tổng quan gần đây~[1,2,3] đều nhấn mạnh: \emph{coastal altimetry remains a scientific and technical challenge}.

\subsection{Vệ tinh không cung cấp chuỗi thời gian liên tục}

Điểm then chốt với \emph{forecasting} là temporal sampling. Hầu hết vệ tinh altimetry có chu kỳ quay lại (revisit time) rất dài:

\begin{table}[ht]
  \centering
  \caption{So sánh temporal sampling giữa trạm và vệ tinh.}
  \label{tab:revisit}
  \begin{tabular}{@{}lcc@{}}
    \toprule
    \textbf{Nguồn dữ liệu} & \textbf{Tần suất} & \textbf{Phù hợp forecasting?} \\
    \midrule
    Tide gauge cố định & Mỗi 6 phút (liên tục) & \textbf{Có} --- capture mọi tần số \\
    Jason-3 / Sentinel-6 & $\sim$10 ngày & Không --- miss toàn bộ diurnal/semi-diurnal \\
    SWOT & $\sim$21 ngày & Không --- dù spatial coverage đẹp \\
    \bottomrule
  \end{tabular}
\end{table}

Vệ tinh \textbf{không thể thay thế} chuỗi đo liên tục tại bờ cho bài toán động lực thời gian nhanh (tidal, surge, seiches).

\subsection{Tide gauges vẫn là nguồn tham chiếu tốt nhất cho sea level ven bờ}

Các tổng quan~[4] vẫn kết luận: \emph{``in-situ tide gauges remains the best approach for long-term coastal sea-level study''}. Altimetry phù hợp hơn ở quy mô regional/global và cần được tích hợp (fusion) với dữ liệu in-situ.

\subsection{Bài toán thực sự: Forecast full field từ cực ít quan trắc cố định}

Bài toán của chúng tôi không chỉ là ``có dữ liệu hay không'' mà là: \textbf{``Có thể forecast toàn bộ trường mực nước từ cực ít trạm quan trắc cố định hay không?''} Điều này vẫn là câu hỏi nghiên cứu mở vì:
\begin{itemize}[nosep]
  \item Trạm bờ rẻ hơn và cung cấp dữ liệu liên tục cho vận hành địa phương.
  \item Nhiều bài toán ven bờ cần latency thấp (cảnh báo sớm sóng triều, storm surge).
  \item Trong thực tế, giải pháp tốt nhất là \emph{fusion} giữa gauges, models, và satellite.
\end{itemize}

\begin{tcolorbox}[colback=green!5, colframe=green!50!black, title={Positioning Statement}]
We study how to forecast coastal/full tidal fields from sparse fixed boundary observations---a setting that remains practically important because coastal in-situ sensors provide continuous local measurements while satellite observations remain intermittent and more challenging in the nearshore zone.
\end{tcolorbox}

% ══════════════════════════════════════════════════════════════
\section{ForecastDeepONet: Kiến trúc và Kết quả}

\subsection{Kiến trúc Causal LSTM Branch}

\begin{table}[ht]
  \centering
  \caption{So sánh kiến trúc Boundary-DeepONet (tuần 7--8) và ForecastDeepONet (tuần 10).}
  \label{tab:arch_compare}
  \begin{tabular}{@{}lcc@{}}
    \toprule
    \textbf{Thành phần} & \textbf{Boundary-DeepONet} & \textbf{ForecastDeepONet} \\
    \midrule
    Branch encoder & MLP (non-causal) & \textbf{LSTM (causal)} \\
    Branch input & $K \times N_t$ flattened & $(T_\mathrm{obs}, K)$ sequential \\
    Temporality & Nhìn toàn bộ $t \in [0,T]$ & Chỉ nhìn $t \leq T_\mathrm{obs}$ \\
    Trunk encoding & Raw $(x,y,t)$, 3D & Fourier PE, 51D \\
    Latent dim $p$ & 128 & 256 \\
    Parameters & $\sim$4,7M & $\sim$1,1M \\
    \midrule
    Task & Reconstruction ($\forall t$) & \textbf{Forecasting ($t > T_\mathrm{obs}$)} \\
    \bottomrule
  \end{tabular}
\end{table}

LSTM xử lý tuần tự sensor time-series từ $t=0$ đến $T_\mathrm{obs}$, xuất hidden state cuối cùng $\mathbf{h}_{T_\mathrm{obs}}$ làm đại diện cho ``trạng thái biển hiện tại''. Đảm bảo \textbf{strict causality}---không bao giờ nhìn trộm tương lai.

\subsection{Training Strategy: Random Horizon Sampling}

Mỗi training step, $T_\mathrm{obs}$ được sample ngẫu nhiên từ khoảng $[T_\mathrm{min}, T_\mathrm{max}]$. Label chỉ lấy từ các timestep $t > T_\mathrm{obs}$. Chiến lược này buộc model phải hoạt động tốt ở \emph{mọi} độ dài quan sát, thay vì overfit vào một horizon cụ thể.

\subsection{Kết quả trên PDEBench SWE 2D}

\subsubsection{Setup}
\begin{itemize}[nosep]
  \item Dataset: PDEBench 2D SWE (Radial Dam Break), 1.000 simulations, grid $128 \times 128$, 101 timesteps.
  \item Sensors: 16 điểm trên 4 cạnh biên (0,1\% diện tích grid).
  \item Train/Val: 800/200 simulations.
  \item Training: 100 epochs, 174 giây, AdamW ($\alpha=10^{-3}$, weight decay $10^{-4}$).
\end{itemize}

\subsubsection{So sánh với Non-Causal Baselines}

\begin{table}[ht]
  \centering
  \caption{Kết quả trên PDEBench 2D SWE. ForecastDeepONet (causal) so với các baseline non-causal.}
  \label{tab:pdebench_results}
  \begin{tabular}{@{}llcc@{}}
    \toprule
    \textbf{Model} & \textbf{Observation} & \textbf{Causality} & \textbf{Rel L2 (\%)} $\downarrow$ \\
    \midrule
    U-Net & Full $128 \times 128$ & Non-causal & 12,80 \\
    FNO-2D & Full $128 \times 128$ & Non-causal & 5,13 \\
    DeepONet (full-state) & Complete IC & Non-causal & 3,74 \\
    Partial-DeepONet & 16 interior sensors & Non-causal & 3,95 \\
    Boundary-DeepONet & 16 boundary sensors & Non-causal & 5,25 \\
    \midrule
    \textbf{ForecastDeepONet} & \textbf{16 boundary sensors} & \textbf{Causal} & \textbf{4,46} \\
    \bottomrule
  \end{tabular}
\end{table}

\subsubsection{Phân tích theo Forecast Horizon}

\begin{table}[ht]
  \centering
  \caption{ForecastDeepONet: Rel L2 theo số timestep quan sát (observation window).}
  \label{tab:horizon}
  \begin{tabular}{@{}cccc@{}}
    \toprule
    $T_\mathrm{obs}$ (quan sát) & $T_\mathrm{future}$ (dự báo) & Rel L2 (\%) & Baseline \\
    \midrule
    20 timesteps & 81 future & 4,72 & 9,80 \\
    40 timesteps & 61 future & 4,43 & 9,52 \\
    60 timesteps & 41 future & 4,33 & 9,82 \\
    80 timesteps & 21 future & 4,46 & 10,23 \\
    \midrule
    \textbf{Trung bình} & & \textbf{4,46} & $\sim$10 \\
    \bottomrule
  \end{tabular}
\end{table}

\subsubsection{Kiểm tra Overfitting}

\begin{table}[ht]
  \centering
  \caption{Train vs.\ Validation Rel L2 --- không có dấu hiệu overfitting.}
  \label{tab:overfit}
  \begin{tabular}{@{}lc@{}}
    \toprule
    & Rel L2 (\%) \\
    \midrule
    Train & 4,38 \\
    Validation & 4,36 \\
    Gap & $-$0,02 \\
    \bottomrule
  \end{tabular}
\end{table}

\subsection{Nhận xét chính}

\begin{enumerate}[nosep]
  \item \textbf{Counter-intuitive:} Causal ForecastDeepONet (4,46\%) \emph{thắng} non-causal Boundary-DeepONet (5,25\%) dù có ít thông tin hơn. LSTM trích xuất temporal dynamics tốt hơn MLP flatten.
  \item \textbf{Không forecast degradation:} Sai số ổn định 4,33--4,72\% qua mọi horizon (21--81 future steps).
  \item \textbf{Không overfit:} Gap train-val = $-$0,02\%. Model 1,1M params vs.\ $\sim$400k samples/epoch.
  \item \textbf{Cực nhanh:} Converge sau 4 epochs, training chỉ 174 giây.
\end{enumerate}


% ══════════════════════════════════════════════════════════════
\section{Tái cấu trúc Paper}

\subsection{Logic mới}

Dựa trên kết quả tuần 10, paper được tái cấu trúc hoàn toàn:

\begin{table}[ht]
  \centering
  \caption{So sánh logic paper cũ và mới.}
  \label{tab:paper_logic}
  \begin{tabular}{@{}lll@{}}
    \toprule
    & \textbf{Logic cũ (tuần 7--9)} & \textbf{Logic mới (tuần 10)} \\
    \midrule
    Bài toán trung tâm & Reconstruction & \textbf{Causal Forecasting} \\
    Đóng góp chính & Boundary-DeepONet & \textbf{ForecastDeepONet (LSTM)} \\
    Benchmark chính & PDEBench (reconstruct) & \textbf{PDEBench (forecast)} \\
    Gulf of Tonkin & Main story & \textbf{Supporting validation} \\
    Novelty claim & Sparse $\to$ field & \textbf{Causal < Non-causal in info, but wins} \\
    \bottomrule
  \end{tabular}
\end{table}

\subsection{Title mới}

\begin{center}
\textit{``Causal Full-Field Forecasting from Sparse Boundary Sensors via Deep Operator Networks''}
\end{center}

\subsection{Trạng thái viết Paper}

\begin{table}[ht]
  \centering
  \caption{Trạng thái các section trong paper sau tuần 10.}
  \label{tab:paper_status}
  \begin{tabular}{@{}lcc@{}}
    \toprule
    \textbf{Section} & \textbf{Trạng thái} & \textbf{Ghi chú} \\
    \midrule
    Title & $\checkmark$ & Đã đổi sang Causal Forecasting \\
    Abstract & TODO & Cần kết quả ablation \\
    Introduction & TODO & Reasoning đã soạn xong (xem \S2) \\
    Related Work & Cập nhật & Thêm coastal altimetry references \\
    Methodology & $\checkmark$ & ForecastDeepONet + Fourier PE \\
    Experiments (PDEBench) & $\checkmark$ & Causal 4,46\% vs.\ non-causal 5,25\% \\
    Gulf of Tonkin (supporting) & $\checkmark$ & 14,8\% Fourier DeepONet \\
    Ablation Studies & TODO & Xem kế hoạch \S5 \\
    Conclusion & TODO & \\
    \bottomrule
  \end{tabular}
\end{table}


% ══════════════════════════════════════════════════════════════
\section{Kế Hoạch Implement Tiếp Theo}

\subsection{Ablation Studies (Ưu tiên cao nhất)}

Để paper đủ lực cho workshop submission, cần bổ sung ablation trên 3 trục chính:

\begin{table}[ht]
  \centering
  \caption{Ma trận ablation dự kiến.}
  \label{tab:ablation_plan}
  \begin{tabular}{@{}llll@{}}
    \toprule
    \textbf{Trục} & \textbf{Configurations} & \textbf{Giả thuyết} & \textbf{Thời gian} \\
    \midrule
    \multicolumn{4}{@{}l}{\textit{Trục quan trắc (observation):}} \\
    Số sensors $K$ & 4, 8, 16, 32, 64 & Diminishing returns $> K=16$ & 1 ngày \\
    Vị trí sensor & Boundary / interior / random & Boundary is sufficient & 0,5 ngày \\
    Sensor noise & $\sigma = 0, 0.01, 0.05, 0.1$ & Robustness & 0,5 ngày \\
    Sensor dropout & 10\%, 30\%, 50\% missing & Graceful degradation & 0,5 ngày \\
    \midrule
    \multicolumn{4}{@{}l}{\textit{Trục kiến trúc (architecture):}} \\
    Branch encoder & LSTM / GRU / TCN / MLP & LSTM tốt nhất & 1 ngày \\
    Fourier freqs $L$ & 0, 4, 8, 16 & Optimal $L=8$ & 0,5 ngày \\
    Latent dim $p$ & 64, 128, 256, 512 & Sweet spot at 256 & 0,5 ngày \\
    \midrule
    \multicolumn{4}{@{}l}{\textit{Trục forecast (temporal):}} \\
    Forecast horizon & $T_\mathrm{obs}/T \in \{0.1, \ldots, 0.8\}$ & Stable (đã verify sơ bộ) & 0,5 ngày \\
    \bottomrule
  \end{tabular}
\end{table}

\subsection{Benchmark bổ sung: OceanBench / GLORYS-CMEMS}

Để tăng sức thuyết phục, dự kiến thêm 1 benchmark trên dữ liệu đại dương mở:
\begin{itemize}[nosep]
  \item \textbf{Ưu tiên 1:} OceanBench (nếu có SSH reanalysis phù hợp).
  \item \textbf{Ưu tiên 2:} GLORYS12 từ Copernicus Marine Service --- dữ liệu reanalysis toàn cầu, có SSH ở độ phân giải 1/12°.
  \item \textbf{Mục tiêu:} Chứng minh ForecastDeepONet transfers từ SWE idealized sang ocean dynamics thực.
\end{itemize}

\subsection{Visualization cho Paper}

\begin{itemize}[nosep]
  \item Heatmap: GT vs.\ Pred vs.\ Error cho 4 timesteps (PDEBench forecasting).
  \item Ablation curves: Rel L2 vs.\ $K$ sensors, Rel L2 vs.\ horizon, Rel L2 vs.\ noise $\sigma$.
  \item Architecture diagram: ForecastDeepONet (LSTM $\to$ dot product $\leftarrow$ Fourier trunk).
\end{itemize}


% ══════════════════════════════════════════════════════════════
\section{Bảng Tổng Hợp Kết Quả Toàn Bộ Nghiên Cứu}

\begin{table}[ht]
  \centering
  \caption{Tổng hợp kết quả từ Tuần 3 đến Tuần 10.}
  \label{tab:summary_all}
  \begin{tabular}{@{}lcccc@{}}
    \toprule
    \textbf{Mô hình} & \textbf{Dataset} & \textbf{Rel L2} & \textbf{Causal?} & \textbf{Tuần} \\
    \midrule
    \multicolumn{5}{@{}l}{\textit{PDEBench SWE 2D (Benchmark):}} \\
    FNO-2D (external) & PDEBench & 5,13\% & --- & 7--8 \\
    DeepONet Full-State & PDEBench & 3,74\% & No & 7--8 \\
    Boundary-DeepONet & PDEBench & 5,25\% & No & 7--8 \\
    \textbf{ForecastDeepONet} & \textbf{PDEBench} & \textbf{4,46\%} & \textbf{Yes} & \textbf{10} \\
    \midrule
    \multicolumn{5}{@{}l}{\textit{Vịnh Bắc Bộ (Case Study):}} \\
    DeepONet (setup gốc) & Vịnh Bắc Bộ & 93,6\% & No & 7--8 \\
    Fourier DeepONet & Vịnh Bắc Bộ & 14,8\% & No & 9 \\
    \bottomrule
  \end{tabular}
\end{table}


% ══════════════════════════════════════════════════════════════
\section*{Tài Liệu Tham Khảo}

\begin{enumerate}[label={[\arabic*]}, nosep]
  \item Vignudelli, S. et al. (2025). Advances and challenges in coastal altimetry. \textit{Frontiers in Marine Science}.
        \url{https://www.frontiersin.org/articles/10.3389/fmars.2025.1592765}
  \item Cipollini, P. et al. (2025). Coastal altimetry---past, present and future. \textit{Ocean Science}, 21, 133--152.
        \url{https://os.copernicus.org/articles/21/133/2025/}
  \item Restano, M. et al. (2025). SAR altimetry for coastal applications. \textit{Int.\ J.\ Applied Earth Obs.}, 106296.
  \item Cazenave, A. et al. (2021). Sea level along the world's coastlines can be measured by a network of tide gauges. \textit{Ocean \& Coastal Management}, 198, 105374.
  \item Su, X. et al. (2025). SWOT temporal sampling limitations for inland waters. \textit{AGU Water Resources Research}.
  \item Lu, L. et al. (2021). Learning nonlinear operators via DeepONet. \textit{Nature Machine Intelligence}, 3, 218--229.
  \item Li, Z. et al. (2021). Fourier Neural Operator for Parametric PDEs. \textit{ICLR 2021}.
  \item Takamoto, M. et al. (2022). PDEBench. \textit{NeurIPS 2022}.
  \item Mildenhall, B. et al. (2020). NeRF. \textit{ECCV 2020}.
\end{enumerate}

\end{document}
